\section{Agentes Inteligentes, MCP, LangGraph e Sistemas Autônomos}

\textbf{OBJETIVO GERAL:}

Desenvolver competências para projetar, desenvolver, integrar, implantar e monitorar agentes inteligentes baseados em Grandes Modelos de Linguagem (LLMs), utilizando arquiteturas modernas de sistemas autônomos, Model Context Protocol (MCP), LangGraph, CrewAI e outros frameworks para construção de aplicações inteligentes capazes de planejar, utilizar ferramentas, gerenciar memória, colaborar com outros agentes e executar tarefas complexas de forma segura, escalável e observável.

\textbf{CARGA HORÁRIA:} 160 horas

\textbf{FUNÇÃO:}

Desenvolver sistemas computacionais baseados em Inteligência Artificial, Análise Preditiva e Ecossistemas IoT, atendendo às normas de qualidade corporativas, robustez metodológica, governança de dados e arquitetura segura.

\textbf{SUBFUNÇÕES:}

\begin{itemize}
\item Projetar arquiteturas de agentes inteligentes.
\item Desenvolver sistemas autônomos baseados em LLMs.
\item Integrar ferramentas externas utilizando MCP.
\item Desenvolver aplicações utilizando LangGraph e CrewAI.
\item Implementar sistemas multiagentes.
\item Gerenciar memória, contexto e planejamento de agentes.
\item Desenvolver workflows inteligentes.
\item Implantar agentes em ambientes produtivos.
\end{itemize}

\textbf{PADRÕES DE DESEMPENHO:}

\begin{itemize}
\item Identificar arquiteturas modernas de agentes inteligentes.
\item Desenvolver agentes utilizando LangGraph.
\item Desenvolver agentes utilizando CrewAI.
\item Integrar ferramentas através do Model Context Protocol.
\item Desenvolver sistemas multiagentes.
\item Implementar memória de curto e longo prazo.
\item Desenvolver workflows autônomos.
\item Aplicar mecanismos de observabilidade.
\item Implementar estratégias de segurança para agentes inteligentes.
\item Documentar arquiteturas desenvolvidas.
\end{itemize}

\textbf{CAPACIDADES TÉCNICAS:}

\textbf{Arquiteturas de Agentes}

\begin{itemize}
\item Interpretar arquiteturas modernas de agentes inteligentes.
\item Projetar agentes baseados em objetivos.
\item Desenvolver agentes reativos e deliberativos.
\item Projetar arquiteturas para sistemas autônomos.
\end{itemize}

\textbf{Frameworks para Agentes}

\begin{itemize}
\item Desenvolver aplicações utilizando LangGraph.
\item Desenvolver aplicações utilizando CrewAI.
\item Utilizar OpenAI Agents SDK.
\item Utilizar PydanticAI e Semantic Kernel.
\end{itemize}

\textbf{Model Context Protocol}

\begin{itemize}
\item Desenvolver aplicações utilizando MCP.
\item Integrar ferramentas externas.
\item Integrar recursos utilizando MCP Servers.
\item Desenvolver aplicações interoperáveis.
\end{itemize}

\textbf{Sistemas Multiagentes}

\begin{itemize}
\item Desenvolver arquiteturas multiagentes.
\item Orquestrar comunicação entre agentes.
\item Implementar agentes especialistas.
\item Desenvolver sistemas cooperativos.
\end{itemize}

\textbf{Planejamento e Memória}

\begin{itemize}
\item Implementar memória de curto prazo.
\item Implementar memória de longo prazo.
\item Desenvolver agentes com planejamento.
\item Aplicar estratégias de raciocínio estruturado.
\end{itemize}

\textbf{Observabilidade e Segurança}

\begin{itemize}
\item Monitorar execução dos agentes.
\item Implementar rastreabilidade.
\item Aplicar mecanismos de segurança.
\item Validar comportamento dos agentes inteligentes.
\end{itemize}

\textbf{CONHECIMENTOS:}

\textbf{Fundamentos de Agentes Inteligentes}

\begin{itemize}
\item Agentes Inteligentes
\item Agentic AI
\item Sistemas Autônomos
\item Arquiteturas Cognitivas
\item Ciclo de Vida dos Agentes
\end{itemize}

\textbf{Arquiteturas de Agentes}

\begin{itemize}
\item Agentes Reativos
\item Agentes Deliberativos
\item Agentes Baseados em Objetivos
\item Agentes Baseados em Utilidade
\item Agentes Cognitivos
\item Arquiteturas Híbridas
\end{itemize}

\textbf{Frameworks}

\begin{itemize}
\item LangGraph
\item CrewAI
\item OpenAI Agents SDK
\item Semantic Kernel
\item PydanticAI
\item AutoGen
\end{itemize}

\textbf{Model Context Protocol}

\begin{itemize}
\item Model Context Protocol
\item MCP Hosts
\item MCP Clients
\item MCP Servers
\item Resources
\item Prompts
\item Tools
\item Sampling
\item Roots
\item Stdio
\item Server Sent Events (SSE)
\item Streamable HTTP
\end{itemize}
